\section{Results and Discussion}
\label{sec:results}

\subsection{RQ1: How Effectively Do Carbon Finance Mechanisms Incentivize Investment and Improve Commercial Viability?}
\label{subsec:rq1}
Effectiveness was assessed through a combination of quantitative data (Section~\ref{subsec:doc-analysis}) and stakeholder perceptions, particularly from developers (Section~\ref{subsec:interviews}). This dual approach allowed for evaluation of how carbon market integration influences financial and operational outcomes in rural solar mini-grids in Kenya and Nigeria. Evidence is presented across these dimensions, highlighting both observed patterns and participant insights.

\subsubsection{Developer Perceptions of Carbon Market Viability}
Table~\ref{tab:awareness} details all four interviewed developers demonstrated awareness of carbon finance despite direct engagement remaining limited.

\begin{table}[h]
	\centering
	\resizebox{\textwidth}{!}{%
	\begin{tabular}{l l l l >{\raggedright\arraybackslash}p{4cm}}
		\hline
		\textbf{Developer} & \textbf{\makecell[lt]{Prior \\ awareness}} & \textbf{\makecell[lt]{Ever \\ considered}} & \textbf{\makecell[lt]{Current \\ status}} & \textbf{\makecell[lt]{Barrier/enabler \\ cited}} \\
		\hline
		\textbf{DEV-NG-01} & Yes, partial & Yes & Not pursuing & Complexity, high transaction costs, small project size \\
		\hline
		\textbf{DEV-NG-02} & Yes & Yes & Attempted & Low emission volumes, high costs, insuffiecient revenue \\
		\hline
		\textbf{DEV-KE-01} & Yes & No & Pursuing & Methodology challenges, high consultatn costs, small portfolio \\
		\hline
		\textbf{DEV-KE-02} & Yes & Yes & Pursuing & Market not accommodating, taxation, verification challenges \\
		\hline
	\end{tabular}}%
	\caption{Developer Awareness and Consideration of Carbon Finance}
	\label{tab:awareness}
\end{table}

Despite recognizing carbon revenue potential, none had successfully integrated carbon finance, with progress consistently constrained by project scale. Smaller mini-grids generate insufficient emissions reductions to justify transaction and monitoring costs, making carbon finance economically unviable. Several developers have opted for Renewable Energy Certificates (RECs) as simpler alternatives. DEV-NG-02 noted, ``Emissions to be avoided were relatively small, hence the transaction could not go through. It requires megawatts of energy projects.''

Expert perspectives reinforced structural barriers: "The cost of entry is high, and fragmentation is a big issue...methodologies weren't really designed with them in mind" (EXP-01). This aligns with Michaelowa et al.'s \cite{michaelowa-catalysing} finding that SSA was "essentially bypassed by the CDM," with only 22\% of single CDM projects and 16\% of PoA component activities issuing credits, a pattern this research confirms persists in voluntary markets.

\subsubsection{Transaction Costs and Certification Barriers}
Table~\ref{tab:costs} analyzes Gold Standard and Verra fee schedules, combined with developer-reported expenses, and reveals a multi-component cost structure creating significant cost barriers for mini-grid carbon certification.

\begin{table}
	\centering
	\resizebox{\textwidth}{!}{%
	\begin{tabular}{l l l l l}
		\hline
		\textbf{Cost Component} & \textbf{<200~kWp} & \textbf{200--500~kWp} & \textbf{>500~kWp} & \textbf{Source} \\
		\hline
		\textbf{Baseline study} & 4--5 & 10--20 & 20--40 & EXP-01, EXP-02 \\
		\hline
		\textbf{PoA registration} & 3.25--4.5 & 3.75--4 & 4--6 & \cite{goldStandardFees,verraFees} \\
		\hline
		\textbf{MRV system setup} & 15--25 & 35--40 & 40--60 & \cite{goldStandardFees,verraFees} \\
		\hline
		\textbf{15 annual verifications} & 18 & 30 & 30 & \cite{goldStandardFees,verraFees} \\
		\hline
		\textbf{Technical assistance} & 10--20 & 20--35 & 35--60 & DEV-KE-01, DEV-NG-02 \\
		\hline
		\textbf{Total} & 53--83 & 95--140 & 140--215 & Calculated totals \\
		\hline
	\end{tabular}}%
	\caption{\textbf{Carbon Certification Transaction Costs in thousands of USD.} Cost estimates compiled from carbon registry fee schedules,~\cite{goldStandardFees,verraFees} consultant expert interviews (EXP-01; EXP-02), and developer-reported project expenses (DEV-KE-01; DEV-NG-02). Verification costs calculated over typical 15-year project crediting period.}
	\label{tab:costs}
\end{table}

To assess the disproportionate burden of transaction costs on different project sizes, Table~\ref{tab:pct-costs} breaks down certification costs against typical mini-grid capital expenditure. CAPEX costs assume the average capital expenditure for mini-grids in SSA is approximately 6,824~USD/kWp.~\cite{amda-bam}

\begin{table}
	\centering
	\begin{tabular}{l r r r}
		\hline
		\textbf{\makecell[lt]{Project \\ Capacity}} & \textbf{\makecell[lt]{CAPEX \\ (kUSD)}} & \textbf{\makecell[lt]{Certification \\ Costs (kUSD)}} & \textbf{\% of CAPEX} \\
		\hline
		\textbf{100~kWp} & 682.4 & 68 & 10\% \\
		\hline
		\textbf{250~kWp} & 1,706 & 117.5 & 6.9\% \\
		\hline
		\textbf{500~kWp} & 3,412 & 177.5 & 5.2\% \\
		\hline
	\end{tabular}
	\caption{Transaction Costs as Percentage of Project CAPEX}
	\label{tab:pct-costs}
\end{table}

This analysis demonstrates that smaller projects face a higher relative burden of certification costs, underscoring the disproportionate impact of transaction costs on small-scale mini-grid developers. For mini-grids under 200~kWp, certification costs represent approximately 10\% of total capital expenditure, which developers consistently identified as a significant constraint on project feasibility. As one small developer noted, ``It requires megawatts of energy projects'' (DEV-NG-02).

Larger developers acknowledged certification costs but emphasized that operational and regulatory barriers often posed more substantial challenges: ``Mini-grids just aren't seen as a viable asset compared to cookstoves...aggregation would help, but it's not easy to pull off'' (EXP-01).

These perspectives indicate that, although transaction costs constitute a meaningful financial barrier, they interact with broader challenges, including project fragmentation, aggregation difficulties, and methodological complexities, that collectively limit the participation of smaller mini-grid developers in carbon markets.

\subsubsection{Revenue Projections}
Small-scale rural mini-grids face significant constraints generating carbon revenue due to limited scale and dispersed nature. As EXP-01 observed: ``Mini-grids just aren't seen as a viable asset.... Intermediaries and registries don't prioritize them because the returns are smaller.'' Current voluntary carbon market pricing of approximately 1--5 USD per \tcooe\ (EXP-02) renders transaction costs systematically greater than potential revenues for sub-200-kWp projects.

While PoA mechanisms theoretically allow small projects to pool credits across jurisdictions, achieving aggregation requires navigating complex registry requirements, demonstrating additionality, and coordinating operations, introducing substantial operational and financial challenges. As POL-02 noted, ``Aggregation is possible using programs of activities, but scalability is challenging because projects take years to join programs, and individual mini-grid volumes are low.'' Post-Paris governance fragmentation, with its competing standards, divergent accounting rules, and uncertain host government authorisation under Article 6, has further complicated the revenue outlook.~\cite{ahonen-governance}

These structural limitations underscore that small rural mini-grids cannot rely on carbon markets as meaningful revenue sources without aggregation mechanisms and supportive policy frameworks that current market structures do not produce organically.

\subsubsection{Summary of RQ1 Findings}
Carbon finance functions as a supplementary rather than transformative factor in mini-grid investment decisions. No developer indicated carbon revenue was decisive; instead, it is considered a ``nice-to-have'' addition to projects pursued independently of carbon finance opportunities. Transaction costs systematically exceed potential revenues for sub-megawatt projects, while revenue leakage means even successful credit sales deliver only a fraction of projected income to developers.~\cite{power-africa-credits}

While additionality is theoretically easier to demonstrate for mini-grids (i.e. projects often would not proceed without additional revenue), existing carbon methodologies were designed for larger-scale installations and require adaptation. Market realities of demand, pricing, and risk perception further constrain investor attention, with buyers not viewing mini-grids as strong asset classes compared to cookstoves or REDD+ projects.

Carbon finance can provide supplemental income but is insufficient as a primary investment incentive for mini-grids. Its impact on commercial viability depends on: (1) project aggregation to achieve economies of scale, (2) methodological adaptations to reduce transaction costs, (3) supportive policy frameworks to address regulatory barriers, (4) cross-border coordination enabling pooled credit volumes, and (5) resolution of post-Paris authorisation requirements under Article 6. Without these mechanisms, carbon finance acts as a secondary rather than standalone financial lever.

\subsection{RQ2: What Conditions Are Necessary for Successful Carbon Market Integration at Scale?}
\label{subsec:rq2}

\subsubsection{Technical Conditions}
MRV requirements present significant barriers for small-scale projects.~\cite{michaelowa2005} As EXP-01 noted: ``MRV is expensive, and the administrative burden is high. For small projects it just eats too much of the value.'' Many developers lack infrastructure and expertise for continuous emissions monitoring, while certifying bodies face knowledge gaps regarding mini-grid verification. Methodologies often consider embodied emissions from imported solar equipment, requiring supply chain tracking capacity that small operators lack (REG-01).

Establishing accurate baselines also proved challenging: ``Grid emission factors, were previously inflating the final emission reduction, so we have to go back and figure out how this reflects in different countries'' (REG-01).

Overcoming scale barriers requires aggregation, yet market mechanisms have failed to produce it organically. While Programme of Activities (PoA) mechanisms theoretically permit incremental project addition, they involve added complexity and scrutiny (REG-01). EXP-01 noted intermediaries focus on larger asset classes: ``If someone could aggregate projects under one umbrella and make it big enough, then maybe intermediaries would pay attention...but right now, they're focused on easier asset classes.''

The atmosfair-MiniGridCo PoA represents the sector's only significant success, succeeding due to four prerequisites absent for independent developers: specialized technical capacity, upfront capital absorbing the 35,000--65,000 USD certification costs, patient capital enabling 10+ year commitments, and registry credibility. Independent Nigerian developers (DEV-NG-01, DEV-NG-02) attempted carbon finance but failed.

Limited replication of the MiniGridCo model three years after its establishment is consistent with, though not conclusive proof of, structural market failure. The nascent timeframe means regulatory and awareness barriers may also be contributing factors. However, DEV-NG-01 and DEV-NG-02 both independently attempted and abandoned carbon finance, citing prohibitive costs and institutional barriers rather than awareness or timing, suggesting structural rather than transitional constraints. This is consistent with World Bank evidence that post-Paris carbon aggregation has become ``smaller, and less transparent, than during the Kyoto Protocol period.''~\cite{michaelowa-catalysing}

\subsubsection{Financial Conditions}
Transaction costs constitute the primary integration barrier. DEV-NG-02 reported: ``The volumes were too low and the revenue to be generated did not match up to the cost we would accrue for the registration and monitoring fees.'' Fixed auditor costs of 100,000--300,000 USD create minimum viable scale thresholds unattainable by individual mini-grids (REG-01).

Evidence on carbon revenue significance varies by project scale. DEV-KE-01 projected that carbon revenues could reach ``12--15\% of total revenue, close to 2 million euros a year once accredited'' for their portfolio. However, DEV-NG-02 reported that actual returns may fall below 5\% of total investment, indicating carbon revenue alone cannot salvage unviable projects. Carbon finance serves as a viability enhancer rather than primary solution, with its contribution depending heavily on project scale, aggregation success, and market pricing. Fundamental mismatches between investor expectations and mini-grid realities further constrain viability: ``Investors also want quick returns, but mini-grids are not a fast-moving business'' (DEV-NG-02).

\subsubsection{Regulatory and Policy Conditions}
POL-02 identified profound regulatory ambiguities in Kenya: ``There's no clear framework or harmonized registry for carbon credits.... Licensing for mini-grids is unclear, including whether developers should be regulated by the Central Bank or other agencies.'' This fragmentation creates uncertainty regarding project eligibility, credit ownership, and operational requirements.

Nigeria presents contrasting approaches. For projects under 1~MWp, ``you get a letter and start operations'' (DEV-NG-02), yet licensing pathway choice has substantive implications: ``If you do registration as a mini grid, you don't have this option [to be a generator]. If you do the permitting option, you can.'' POL-02 emphasized that ``power purchase agreements, leasing, and registration processes are ambiguous, especially for projects under 20 years,'' directly affecting carbon projects with 7--10 year crediting periods.

In Kenya, the Energy and Petroleum Regulatory Authority ``(EPRA) holds issuance of tariffs and approvals.... There is pressure to push tariffs down'' (DEV-KE-01), with stark variation across locations. The relationship between carbon revenues and tariff regulation remains undefined. If regulators treat carbon income as general project revenue, they may mandate tariff reductions, effectively capturing carbon revenues. Conversely, if treated as separate environmental benefit streams, developers retain them. POL-02 confirmed this ambiguity exists with no specific provisions. DEV-KE-01 described regulatory responsiveness challenges: ``Slow response, EPRA hasn't replied.... They are not ready and don't understand.''

REG-01 explained a critical barrier: ``Check [National Determined Contributions] (NDCs), say Kenya has claimed clean energy activity, your project is not additional.'' If governments include mini-grid deployment in unconditional NDC commitments, projects automatically fail additionality tests. POL-02 noted Kenya's experience with benefit-sharing requirements affecting Nature-Based Solutions projects. This creates perverse incentives: governments demonstrating climate ambition inadvertently disqualify domestic projects from carbon finance that could help achieve those targets, exemplifying the policy integration gap where energy and carbon market regulations operate in silos.

Post-Paris governance has intensified this challenge. Article 6.2 of the Paris Agreement requires host government authorization for internationally transferred mitigation outcomes (ITMOs), and host countries now ``bear an unprecedented task of ensuring the environmental integrity and robust accounting of ITMOs.''~\cite{ahonen-governance} Many SSA governments lack this capacity; Kenya has approved zero Article 6 projects despite its Climate Change Act (2016, amended 2023) and Carbon Markets Regulations (2024), which fix Corresponding Adjustment prices at USD 4 per \tcooe\ and strictly limiting Letters of Authorisation to preserve its carbon budget.

The consequences of this authorization barrier are severe. KOKO Networks, serving 1.5 million Kenyan households with bioethanol cooking fuel, entered administration in February 2026 after the Kenyan government refused a Letter of Authorisation, barring monetization under Article 6.2. Seven hundred employees lost jobs and 1.5 million households returned to charcoal and kerosene.~\cite{koko-disrupt} While KOKO Networks was not a mini-grid developer, this case directly validates the authorization risks identified in this research: even real, additional, verifiable emission reductions cannot succeed without host government authorization. The case illustrates the ``extreme fragility of business models that rely on the complex and often shifting regulatory landscapes of international carbon markets.''~\cite{koko-econews}

Grid arrival risk, which implies national grid extension reaching mini-grid areas, poses significant concerns for projects with 10--15-year carbon crediting periods. Nigeria's compensation framework pays ``the depreciated value plus one year of revenue'' (DEV-NG-02), failing to capture multi-year losses with no provision for lost carbon revenues, which one developer projected could reach 12--15\% of total revenue for larger portfolios. This creates fundamental incompatibility between carbon market requirements demonstrating permanence and multi-year monitoring commitments, and grid arrival policies that can unilaterally terminate these commitments.

In Kenya, carbon credit sales incur taxation between 25 and 40\%,~\cite{kenya-climate-change-regulations} substantially reducing net benefits. For example, DEV-KE-01's projected 12--15\% revenue contribution would decrease to approximately 9--11\% after taxation. POL-02 observed this ``discourages developers from selling credits locally,'' potentially incentivizing offshore sales. This taxation approach contradicts stated policy objectives. If carbon markets provide additional financing for marginal projects, taxing revenues at standard rates undermines additionality by capturing 25--40\% of the incremental benefits.

\subsubsection{Carbon Market-Specific Conditions}
Current carbon methodologies were not designed for mini-grids. Registries ``use CDM methodologies, they have been doing them for years'' (REG-01), but these originated for larger projects. Scale thresholds create barriers: projects above 60,000~\tcooe\ annually face stricter rules, while most individual mini-grids fall into microscale categories below 10,000~\tcooe. EXP-02 noted geographic aggregation precedents exist in regenerative agriculture but have not been systematically applied to renewable energy mini-grids. REG-01 confirmed ``registries are now designing for new RE and that is currently ongoing, decision made in 2023,'' indicating methodological adaptation is occurring but incomplete. The shift toward context-specific additionality assessment improves accuracy but increases complexity and transaction costs. The retroactive restriction proves particularly problematic: ``many projects have been running since 2011--2013, so they are not valid for carbon credits'' (DEV-KE-01), creating a chicken-and-egg financing problem.

Buyer preferences significantly shape viability. EXP-02 noted corporate buyers ``focus a lot on permanence'' and prefer removal over reduction credits, disadvantaging renewable energy projects. Recent price quotes ``around one U.S. dollar'' per \tcooe\ were described as ``ridiculous.'' Reputational risk aversion favors established project types, but mini-grids lack sufficient track record for rating agency assessment.

EXP-01 identified the fundamental issue: ``SSA will be a net seller, and the global north will be the buyers, that means the demand side decides what's acceptable.'' This power asymmetry means mini-grid viability depends on shifting buyer preferences in developed markets rather than demonstrating value to host country stakeholders.

Multiple respondents identified Renewable Energy Certificates as potentially more viable. REG-01 advised ``RECs are a safer choice, better regulated, less change and rules around them.'' DEV-KE-02 cited the R-REC Standard~\cite{rrecs} with ``methodology in line with UNFCCC guidelines that allows you to convert RECs to carbon credits,'' representing hybrid flexibility. Platform developments integrating ``I-RECs, TIGRs, and GEOs onto the same platform as carbon credits'' could increase liquidity (EXP-02). REG-01 noted D-RECs command premium prices ``because of their use in specialized areas'' like schools and clinics, suggesting co-benefit documentation enhances value.

EXP-02 described settlement innovations addressing transaction risks: ``We hold the credits...but they still remain the seller's, and then the buyer needs to transfer the money...two days for the payment to happen instead of 30.'' Platform neutrality and bank-grade requirements ensure trustworthy infrastructure, while multi-registry integration reduces transaction costs for buyers and sellers who previously needed separate accounts with each registry.

\subsubsection{Summary of RQ2 Findings}
Analysis reveals successful carbon market integration requires simultaneous convergence across five interdependent condition clusters. The atmosfair-MiniGridCo model's singular success alongside 200 other non-participating Nigerian mini-grids reveals a critical finding: aggregation mechanisms do not emerge organically from market forces.

The critical insight is MiniGridCo's limited replication is consistent with structural market inadequacy, corroborated by DEV-NG-01 and DEV-NG-02 citing prohibitive costs and institutional barriers rather than awareness or timing. Aggregation mechanisms do not emerge organically from market forces despite clear economic rationale and available PoA frameworks. This implies that without deliberate institutional intervention creating aggregation infrastructure, carbon finance remains inaccessible to the mini-grid sector regardless of individual developer capabilities or methodological improvements.

Three insights emerge:
\begin{enumerate}
	\item Market forces do not produce necessary aggregation infrastructure. MiniGridCo succeeded because atmosfair possesses institutional endowments, such as DFI backing, technical expertise, patient capital, unavailable to commercial developers, not because PoA mechanisms are accessible.
	\item Transaction cost economics remain prohibitive. Even with aggregation theoretically available, the existence of PoA pathways has not catalyzed adoption because upfront costs and complexity exceed independent developer capacity.
	\item Regulatory fragmentation creates structural incompatibilities that technical solutions alone cannot overcome, multiple policy domains must align simultaneously.
\end{enumerate}

Table~\ref{tab:rq2} synthesizes barriers and required changes, revealing that scale, aggregation, and transaction cost economics represent binding structural constraints rather than solvable technical problems. Addressing only one without the other yields insufficient improvement, while both depend on institutional capacity that current market structures do not produce.

\begin{table}
	\centering
	\resizebox{\textwidth}{!}{%
	\begin{tabular}{l >{\raggedright\arraybackslash}p{4cm} l >{\raggedright\arraybackslash}p{4cm} >{\raggedright\arraybackslash}p{4cm}}
		\hline
		\textbf{\makecell[lt]{Condition \\ domain}} & \textbf{Structural barrier} & \textbf{Priority} & \textbf{Dependencies} & \textbf{Required intervention} \\
		\textbf{\makecell[lt]{Scale and \\ aggregation}} & Individual projects economically unviable; PoA mechanisms exist but markets have not produced aggregation intermediaries despite 200+ mini-grids operating. & Critical & Requires institutional capacity that markets do not produce organically; depends on public or DFI support. & Establish publicly-funded aggregation intermediaries in SSA with mandate to absorb upfront costs and coordinate multi-developer PoAs. Market forces have proven insufficient. \\
		\hline
		\textbf{\makecell[lt]{Transaction \\ cost \\ economics}} & Fixed costs (100--300k USD) exceed small project revenues at 1--5 USD/\tcooe\ pricing. & Critical & Depends on aggregation and improved pricing. & Public subsidy of certification costs until scale achieved; registry fee waivers for aggregated projects; simplified MRV protocols requiring policy mandate. \\
		\hline
		\textbf{\makecell[lt]{Regulatory \\ framework}} & Carbon revenue treatment undefined in tariff regulations; grid arrival compensation excludes lost carbon revenues (potentially 12--15\% of income); 25--40\% taxation undermines additionality. & High & Enables project economics if combined with scale solutions. & Explicit carbon revenue provisions in licensing; tax exemptions or reduced rates. \\
		\hline
		\textbf{\makecell[lt]{Carbon \\ market \\ design}} & Methodologies designed for large projects; buyer preference for removal over reduction credits. & High & Requires registry willingness to adapt; influenced by buyer demand. & Mini-grid-specific methodologies with proportional MRV; REC-carbon hybrid instruments; revision of retroactive timeline restrictions; recognition of additionality in energy-poor contexts. \\
		\hline
		\textbf{\makecell[lt]{Institutional \\ capacity}} & No specialized intermediary serving mini-grid sector; registry access barriers for new entrants. & Medium & Becomes relevant only after critical constraints addressed; market has not produced this capacity organically. & Adoption of a mini-grid-specific standard to codify issuance rules, reduce verification overhead, and enable registry access via technical interoperability. \\
		\hline
	\end{tabular}}%
	\caption{Synthesis of Barriers and Required Interventions for Carbon Market Integration}
	\label{tab:rq2}
\end{table}

\subsection{Research Limitations}
This qualitative study provides empirically grounded insights but has several scope and generalizability constraints.
\begin{itemize}
	\item \textbf{Sample size and selection:} Four developers (4 of 7 contacted) provide depth but limit generalizability, with potential selection bias toward carbon-engaged developers. Absence of DFI perspectives constrains understanding of investment decision-making. Claims about scale conditions require validation through larger quantitative research.
	\item \textbf{Geographic specificity:} Findings reflect Kenya's and Nigeria's regulatory frameworks and are not directly generalizable to different policy environments.
	\item \textbf{Financial data precision:} Transaction cost analysis combines registry schedules, developer reports, and industry estimates without independent audit. Revenue projections involve uncertainty from future prices and demand.
	\item \textbf{Temporal context:} Data collection (June--August 2025) during carbon market transition. Registry methodology revisions incomplete; 1 USD/\tcooe\ pricing reflects research-time conditions.
\end{itemize}

Despite these limitations, multi-stakeholder design and cross-country comparison provide robust empirical evidence where such research remains scarce.